\subsection{Boundary-crossing component constraints}\label{subsec:block_diagonal_boundary_crossing}

\begin{lemma}[Cyclic restrictions of block-diagonal boundary vectors]
    \label{lem:block_diagonal_boundary_cyclic_restriction_sum}
    \lean{MPSTensor.cyclicRestrictₗ_groundSpaceMap_toTensorFromBlocks_blockDiagonal_eq_sum}
    \leanok
    \uses{lem:block_diagonal_boundary_condition_sum}
    Let \(B=\bigoplus_{j=1}^g\mu_jA_j\), and fix block-diagonal boundary
    conditions \(X_j\). For every cyclic restriction \(R_{i,\tau}\) from
    \(N\) sites to \(L\) sites,
    \[
        R_{i,\tau}\!\left(
        \Gamma_N^B\!\left(\bigoplus_{j=1}^gX_j\right)
        \right)
        =
        \sum_{j=1}^g
        R_{i,\tau}\!\left(\Gamma_N^{A_j}(\mu_j^NX_j)\right).
    \]
\end{lemma}

\begin{proof}
    \leanok
    \uses{lem:block_diagonal_boundary_condition_sum}
    Apply Lemma~\ref{lem:block_diagonal_boundary_condition_sum} at length \(N\),
    then use linearity of the cyclic restriction \(R_{i,\tau}\).
\end{proof}

\begin{lemma}[Local direct-sum constraint from block-diagonal boundary conditions]
    \label{lem:block_diagonal_boundary_local_sum_mem}
    \lean{MPSTensor.blockDiagonal_boundary_cyclicRestrict_sum_mem_iSup_groundSpace}
    \leanok
    \uses{lem:block_diagonal_boundary_cyclic_restriction_sum,
        lem:block_sum_local_ground_space}
    Let \(B=\bigoplus_{j=1}^g\mu_jA_j\), with \(\mu_j\ne0\), and suppose
    \[
        \psi=\Gamma_N^B\!\left(\bigoplus_{j=1}^gX_j\right),
        \qquad
        \psi\in\mathcal G_{N,L}(B).
    \]
    Then every local constraint of \(\psi\) gives
    \[
        \sum_{j=1}^g
        R_{i,\tau}\!\left(\Gamma_N^{A_j}(\mu_j^NX_j)\right)
        \in
        \bigvee_{j=1}^gG_L(A_j).
    \]
\end{lemma}

\begin{proof}
    \leanok
    \uses{lem:block_diagonal_boundary_cyclic_restriction_sum,
        lem:block_sum_local_ground_space}
    The hypothesis \(\psi\in\mathcal G_{N,L}(B)\) says that the left-hand side
    before decomposition lies in \(G_L(B)\).  Lemma
    \ref{lem:block_diagonal_boundary_cyclic_restriction_sum} rewrites this local
    vector as the displayed sum, and Lemma~\ref{lem:block_sum_local_ground_space}
    identifies \(G_L(B)\) with \(\bigvee_jG_L(A_j)\).

    % Next source step, not proven by this membership lemma: blockwise
    % extraction
    % $A^j_{i_{m+1}}C^j_{i_1}=D^j_{i_{m+1}}A^j_{i_1}$
    % from \cite[Theorem~2blocks.2]{PerezGarcia2007Matrix}.
\end{proof}

\begin{lemma}[Cyclic windows not crossing the cut]
    \label{lem:block_diagonal_boundary_nonwrapping_component}
    \lean{MPSTensor.blockDiagonal_boundary_cyclicRestrict_component_mem_groundSpace_of_nonwrapping}
    \leanok
    \uses{def:ground_space_parent, rem:cyclic_window_operators}
    Fix block-diagonal boundary conditions $X_j$. If the cyclic window beginning
    at $i$ stays inside the chosen linear interval, $i+L\le N$, then
    for every block $j$,
    \[
        R_{i,\tau}\!\left(\Gamma_N^{A_j}(\mu_j^NX_j)\right)
        \in G_L(A_j).
    \]
\end{lemma}

\begin{proof}
    \leanok
    \uses{def:ground_space_parent, rem:cyclic_window_operators}
    When the window stays inside the chosen linear interval, the boundary matrix
    remains outside the window. Let $A_{\rm left}$ be the product on the sites
    before the window and $A_{\rm right}$ the product on the sites after the
    window. Trace cyclicity
    gives
    \[
        R_{i,\tau}\!\left(\Gamma_N^{A_j}(\mu_j^NX_j)\right)
        =
        \Gamma_L^{A_j}\!\left(A_{\rm right}\,\mu_j^NX_j\,A_{\rm left}\right).
    \]
    Hence the restricted vector lies in $G_L(A_j)$.
    % Crossing-window note, not part of this non-wrapping lemma: after trace
    % rotation, $\mu_j^N X_j$ lies between the two pieces of the interval. The
    % missing block-diagonal boundary-closing comparison is
    % $A^j_{i_{m+1}}C^j_{i_1}=D^j_{i_{m+1}}A^j_{i_1}$ from
    % \cite[Theorem~2blocks.2]{PerezGarcia2007Matrix}.
\end{proof}

\begin{lemma}[Boundary-crossing cyclic intervals from the boundary matrix identity]
    \label{lem:block_diagonal_boundary_crossing_matrix_identity}
    \lean{MPSTensor.blockDiagonal_boundary_cyclicRestrict_component_mem_groundSpace_of_crossing_matrix}
    \leanok
    \uses{def:ground_space_parent, rem:cyclic_window_operators}
    Assume \(0<N\), \(L\le N\), and let the cyclic interval beginning at \(i\)
    cross the boundary cut, so \(N<i+L\). Put \(a=i+L-N\). Suppose that there is
    a matrix \(E\) such that for every word \(\beta\) of length \(a\),
    \[
        \mu_j^N X_j A^j_\beta
        A^j_{\tau_a}\cdots A^j_{\tau_{i-1}}
        =
        A^j_\beta E.
    \]
    Then
    \[
        R_{i,\tau}\!\left(\Gamma_N^{A_j}(\mu_j^NX_j)\right)
        \in G_L(A_j).
    \]
\end{lemma}

\begin{proof}
    \leanok
    \uses{def:ground_space_parent, rem:cyclic_window_operators}
    Write a length-\(L\) word in the boundary-crossing cyclic interval as
    \(\alpha\beta\), where \(\alpha\) is the segment inside
    \(\{i,\ldots,N-1\}\) and \(\beta\) is the segment after the interval wraps
    past the cut, inside \(\{0,\ldots,a-1\}\). The full
    \(N\)-site word appearing in the restriction is then
    \[
        \beta\,\tau_a\cdots\tau_{i-1}\,\alpha .
    \]
    Trace cyclicity gives
    \[
        \operatorname{tr}\!\left(
        A^j_\beta A^j_{\tau_a}\cdots A^j_{\tau_{i-1}}
        A^j_\alpha \mu_j^NX_j\right)
        =
        \operatorname{tr}\!\left(
        A^j_\alpha \mu_j^NX_j A^j_\beta
        A^j_{\tau_a}\cdots A^j_{\tau_{i-1}}\right).
    \]
    The assumed boundary matrix identity changes the last expression to
    \[
        \operatorname{tr}\!\left(A^j_\alpha A^j_\beta E\right),
    \]
    which is the coefficient of \(\Gamma_L^{A_j}(E)\) at the word
    \(\alpha\beta\). Hence the restricted vector lies in \(G_L(A_j)\).
\end{proof}

\begin{lemma}[Last crossing-window component trace form]
    \label{lem:block_diagonal_last_crossing_component_trace_form}
    \lean{MPSTensor.blockDiagonal_boundary_last_cyclicRestrict_component_eq_leftBoundaryComponent}
    \leanok
    \uses{def:ground_space_map, rem:cyclic_window_operators}
    Write \(N=M+1\) and let the local length be \(m+2\le M+1\). Consider the
    cyclic interval beginning at the last site \(M\). For an outside
    configuration \(\tau\), put
    \[
        \rho=\tau_{m+1}\cdots\tau_{M-1},
        \qquad
        C^j_a=A^j_\rho A^j_a(\mu_j^{M+1}X_j).
    \]
    Then the \(j\)-th component of the last crossing-window restriction has
    the left-boundary trace form
    \[
        R_{M,\tau}\!\left(\Gamma_{M+1}^{A_j}(\mu_j^{M+1}X_j)\right)
        =
        \Phi^j,
        \qquad
        \Phi^j(a,w,b)=\operatorname{tr}\!\left(A^j_b C^j_a A^j_w\right).
    \]
\end{lemma}

\begin{proof}
    \leanok
    \uses{def:ground_space_map, rem:cyclic_window_operators}
    The cyclic interval beginning at \(M\) reads the last site first and then
    wraps to the sites \(0,\ldots,m\). Thus, for the local word \(awb\), the
    full word in the trace is
    \[
        w\,b\,\rho\,a .
    \]
    Hence the \(j\)-th coefficient is
    \[
        \operatorname{tr}\!\left(A^j_wA^j_bA^j_\rho A^j_a
        (\mu_j^{M+1}X_j)\right).
    \]
    Trace cyclicity rotates this to
    \[
        \operatorname{tr}\!\left(
        A^j_bA^j_\rho A^j_a(\mu_j^{M+1}X_j)A^j_w\right),
    \]
    which is the displayed left-boundary trace form.
\end{proof}

\begin{lemma}[Last crossing-window trace form]
    \label{lem:block_diagonal_last_crossing_trace_form}
    \lean{MPSTensor.blockDiagonal_boundary_last_cyclicRestrict_sum_eq_leftBoundaryComponents}
    \leanok
    \uses{lem:block_diagonal_last_crossing_component_trace_form}
    Write \(N=M+1\) and let the local length be \(m+2\le M+1\). Consider the
    cyclic interval beginning at the last site \(M\). For an outside
    configuration \(\tau\), put
    \[
        \rho=\tau_{m+1}\cdots\tau_{M-1},
        \qquad
        C^j_a=A^j_\rho A^j_a(\mu_j^{M+1}X_j).
    \]
    Then the block sum of the last crossing-window restrictions has the
    left-boundary trace form
    \[
        \sum_j
        R_{M,\tau}\!\left(\Gamma_{M+1}^{A_j}(\mu_j^{M+1}X_j)\right)
        =
        \sum_j \Phi^j,
        \qquad
        \Phi^j(a,w,b)=\operatorname{tr}\!\left(A^j_b C^j_a A^j_w\right).
    \]
\end{lemma}

\begin{proof}
    \leanok
    \uses{lem:block_diagonal_last_crossing_component_trace_form}
    Sum Lemma~\ref{lem:block_diagonal_last_crossing_component_trace_form} over
    the block index.
\end{proof}

\begin{lemma}[Last crossing-window coefficient comparison]
    \label{lem:block_diagonal_last_crossing_coefficients_from_sum_membership}
    \lean{MPSTensor.blockDiagonal_boundary_last_coefficients_of_sum_mem_iSup}
    \leanok
    \uses{lem:block_diagonal_last_crossing_trace_form,
        lem:blockwise_boundary_identities_from_left_trace_membership}
    In the setting of Lemma~\ref{lem:block_diagonal_last_crossing_trace_form},
    assume that the length-\(m\) simultaneous word tuples span the product
    algebra and that
    \[
        \sum_j
        R_{M,\tau}\!\left(\Gamma_{M+1}^{A_j}(\mu_j^{M+1}X_j)\right)
        \in \bigvee_j G_{m+2}(A^j).
    \]
    Then there are matrices \(E_j\) such that, for every block \(j\) and all
    boundary letters \(a,b\),
    \[
        A^j_b\bigl(A^j_\rho A^j_a(\mu_j^{M+1}X_j)\bigr)
        =
        A^j_bE_jA^j_a .
    \]
\end{lemma}

\begin{proof}
    \leanok
    \uses{lem:block_diagonal_last_crossing_trace_form,
        lem:blockwise_boundary_identities_from_left_trace_membership}
    Lemma~\ref{lem:block_diagonal_last_crossing_trace_form} writes the local
    block sum in the form
    \[
        \Phi(a,w,b)=
        \sum_j\operatorname{tr}\!\left(A^j_bC^j_aA^j_w\right),
        \qquad
        C^j_a=A^j_\rho A^j_a(\mu_j^{M+1}X_j).
    \]
    Applying
    Lemma~\ref{lem:blockwise_boundary_identities_from_left_trace_membership}
    gives the displayed blockwise equations.
\end{proof}

\begin{lemma}[Last crossing-window component membership]
    \label{lem:block_diagonal_last_crossing_component_membership_from_sum_membership}
    \lean{MPSTensor.blockDiagonal_boundary_last_cyclicRestrict_component_mem_groundSpace_of_sum_mem_iSup}
    \leanok
    \uses{lem:block_diagonal_last_crossing_component_trace_form,
        lem:block_diagonal_last_crossing_coefficients_from_sum_membership,
        lem:left_boundary_component_mem_ground_space}
    In the setting of Lemma~\ref{lem:block_diagonal_last_crossing_trace_form},
    assume that the length-\(m\) simultaneous word tuples span the product
    algebra and that
    \[
        \sum_j
        R_{M,\tau}\!\left(\Gamma_{M+1}^{A_j}(\mu_j^{M+1}X_j)\right)
        \in \bigvee_j G_{m+2}(A^j).
    \]
    Then, for every block \(j\),
    \[
        R_{M,\tau}\!\left(\Gamma_{M+1}^{A_j}(\mu_j^{M+1}X_j)\right)
        \in G_{m+2}(A^j).
    \]
\end{lemma}

\begin{proof}
    \leanok
    \uses{lem:block_diagonal_last_crossing_component_trace_form,
        lem:block_diagonal_last_crossing_coefficients_from_sum_membership,
        lem:left_boundary_component_mem_ground_space}
    Lemma~\ref{lem:block_diagonal_last_crossing_coefficients_from_sum_membership}
    gives matrices \(E_j\) with
    \[
        A^j_bC^j_a=A^j_bE_jA^j_a.
    \]
    By Lemma~\ref{lem:block_diagonal_last_crossing_component_trace_form}, the
    \(j\)-th cyclic restriction is the left-boundary component determined by
    \(C^j_a\). The displayed identity makes that component equal to
    \(\Gamma_{m+2}^{A_j}(E_j)\), hence it lies in \(G_{m+2}(A^j)\).
\end{proof}

\begin{lemma}[Last crossing-window coefficients from a block-diagonal chain vector]
    \label{lem:block_diagonal_last_crossing_coefficients_from_chain_ground_space}
    \lean{MPSTensor.blockDiagonal_boundary_last_coefficients_of_blockDiagonal_chainGroundSpace}
    \leanok
    \uses{lem:block_diagonal_boundary_local_sum_mem,
        lem:block_diagonal_last_crossing_coefficients_from_sum_membership}
    Let \(B=\bigoplus_j\mu_jA^j\), and assume \(\mu_j\ne0\) for every block
    \(j\). Suppose
    \[
        \psi=\Gamma_{M+1}^{B}\!\left(\bigoplus_jX_j\right),
        \qquad
        \psi\in\mathcal G_{M+1,m+2}(B).
    \]
    Assume also that the length-\(m\) simultaneous word tuples span the product
    algebra. Then, for every outside configuration \(\tau\), with
    \(\rho=\tau_{m+1}\cdots\tau_{M-1}\), there are matrices \(E_j\) such that
    \[
        A^j_b\bigl(A^j_\rho A^j_a(\mu_j^{M+1}X_j)\bigr)
        =
        A^j_bE_jA^j_a
    \]
    for every block \(j\) and all boundary letters \(a,b\).
\end{lemma}

\begin{proof}
    \leanok
    \uses{lem:block_diagonal_boundary_local_sum_mem,
        lem:block_diagonal_last_crossing_coefficients_from_sum_membership}
    The local constraint of \(\psi\) at the cyclic interval beginning at \(M\)
    gives
    \[
        \sum_j
        R_{M,\tau}\!\left(\Gamma_{M+1}^{A_j}(\mu_j^{M+1}X_j)\right)
        \in \bigvee_j G_{m+2}(A^j)
    \]
    by Lemma~\ref{lem:block_diagonal_boundary_local_sum_mem}. Applying
    Lemma~\ref{lem:block_diagonal_last_crossing_coefficients_from_sum_membership}
    gives the displayed equations.
\end{proof}

\begin{lemma}[Last crossing-window component membership from a chain vector]
    \label{lem:block_diagonal_last_crossing_component_membership_from_chain_ground_space}
    \lean{MPSTensor.blockDiagonal_boundary_last_component_mem_groundSpace_of_blockDiagonal_chainGroundSpace}
    \leanok
    \uses{lem:block_diagonal_boundary_local_sum_mem,
        lem:block_diagonal_last_crossing_component_membership_from_sum_membership}
    Let \(B=\bigoplus_j\mu_jA^j\), and assume \(\mu_j\ne0\) for every block
    \(j\). Suppose
    \[
        \psi=\Gamma_{M+1}^{B}\!\left(\bigoplus_jX_j\right),
        \qquad
        \psi\in\mathcal G_{M+1,m+2}(B).
    \]
    If the length-\(m\) simultaneous word tuples span the product algebra, then
    for every outside configuration \(\tau\) and every block \(j\),
    \[
        R_{M,\tau}\!\left(\Gamma_{M+1}^{A_j}(\mu_j^{M+1}X_j)\right)
        \in G_{m+2}(A^j).
    \]
\end{lemma}

\begin{proof}
    \leanok
    \uses{lem:block_diagonal_boundary_local_sum_mem,
        lem:block_diagonal_last_crossing_component_membership_from_sum_membership}
    Lemma~\ref{lem:block_diagonal_boundary_local_sum_mem} applied to the
    window beginning at the last site gives
    \[
        \sum_j
        R_{M,\tau}\!\left(\Gamma_{M+1}^{A_j}(\mu_j^{M+1}X_j)\right)
        \in \bigvee_j G_{m+2}(A^j).
    \]
    Lemma~\ref{lem:block_diagonal_last_crossing_component_membership_from_sum_membership}
    then applies to each block component.
\end{proof}

\begin{lemma}[Boundary-crossing intervals from a right boundary-matrix identity]
    \label{lem:block_diagonal_boundary_crossing_right_boundary_matrix}
    \lean{MPSTensor.blockDiagonal_boundary_cyclicRestrict_component_mem_groundSpace_of_boundary_identity}
    \leanok
    \uses{lem:block_diagonal_boundary_crossing_matrix_identity}
    Assume \(0<N\), \(L\le N\), and let the cyclic interval beginning at \(i\)
    cross the boundary cut, so \(N<i+L\). Put \(a=i+L-N\). Suppose that there is
    a matrix \(Y\) such that for every word \(\beta\) of length \(a\),
    \[
        \mu_j^N X_j A^j_\beta=A^j_\beta Y .
    \]
    Then
    \[
        R_{i,\tau}\!\left(\Gamma_N^{A_j}(\mu_j^NX_j)\right)
        \in G_L(A_j).
    \]
\end{lemma}

\begin{proof}
    \leanok
    \uses{lem:block_diagonal_boundary_crossing_matrix_identity}
    Let
    \[
        E:=Y A^j_{\tau_a}\cdots A^j_{\tau_{i-1}} .
    \]
    Then, for every word \(\beta\) on the segment \(0,\ldots,a-1\),
    \[
        \mu_j^N X_j A^j_\beta
        A^j_{\tau_a}\cdots A^j_{\tau_{i-1}}
        =
        A^j_\beta E.
    \]
    Lemma~\ref{lem:block_diagonal_boundary_crossing_matrix_identity} gives the
    stated local constraint.
\end{proof}

\begin{lemma}[Boundary-crossing cyclic intervals suffice]
    \label{lem:block_diagonal_boundary_component_crossing_windows_suffice}
    \lean{MPSTensor.blockDiagonal_boundary_component_chainGroundSpace_of_crossing_windows}
    \leanok
    \uses{def:chain_ground_space,
        lem:block_diagonal_boundary_nonwrapping_component}
    Assume \(0<N\) and \(L\le N\). Fix block-diagonal boundary conditions
    \(X_j\). Suppose that, for every block \(j\), every cyclic interval
    beginning at \(i\) with \(N<i+L\), and every configuration \(\tau\),
    \[
        R_{i,\tau}\!\left(\Gamma_N^{A_j}(\mu_j^NX_j)\right)
        \in G_L(A_j).
    \]
    Then, for every block \(j\),
    \[
        \Gamma_N^{A_j}(\mu_j^NX_j)\in \mathcal G_{N,L}(A_j).
    \]
\end{lemma}

\begin{proof}
    \leanok
    \uses{def:chain_ground_space,
        lem:block_diagonal_boundary_nonwrapping_component}
    By definition, membership in \(\mathcal G_{N,L}(A_j)\) is the collection of
    all cyclic length-\(L\) local constraints. For an interval beginning at \(i\),
    either \(i+L\le N\) or \(N<i+L\). In the first case use
    Lemma~\ref{lem:block_diagonal_boundary_nonwrapping_component}. In the second
    case use the displayed hypothesis.
\end{proof}

\begin{lemma}[Right boundary-matrix identities give componentwise periodic constraints]
    \label{lem:block_diagonal_boundary_component_right_boundary_matrices}
    \lean{MPSTensor.blockDiagonal_boundary_component_chainGroundSpace_of_boundary_identities}
    \leanok
    \uses{lem:block_diagonal_boundary_crossing_right_boundary_matrix,
        lem:block_diagonal_boundary_component_crossing_windows_suffice}
    Assume \(0<N\) and \(L\le N\). Fix block-diagonal boundary conditions
    \(X_j\). Suppose that, for every block \(j\), every boundary-crossing cyclic
    interval beginning at \(i\), and every configuration \(\tau\), there is a
    matrix \(Y_{j,i,\tau}\) such that, for every word \(\beta\) of length
    \(a=i+L-N\),
    \[
        \mu_j^N X_j A^j_\beta
        =
        A^j_\beta Y_{j,i,\tau}
    \]
    Then,
    for every block \(j\),
    \[
        \Gamma_N^{A_j}(\mu_j^NX_j)\in \mathcal G_{N,L}(A_j).
    \]
\end{lemma}

\begin{proof}
    \leanok
    \uses{lem:block_diagonal_boundary_crossing_right_boundary_matrix,
        lem:block_diagonal_boundary_component_crossing_windows_suffice}
    Lemma~\ref{lem:block_diagonal_boundary_crossing_right_boundary_matrix}
    turns the displayed boundary identity into the local constraint for each
    boundary-crossing interval.
    Lemma~\ref{lem:block_diagonal_boundary_component_crossing_windows_suffice}
    then yields
    \[
        \Gamma_N^{A_j}(\mu_j^NX_j)\in \mathcal G_{N,L}(A_j),
    \]
    handling both boundary-crossing intervals \(N<i+L\) and non-crossing
    intervals \(i+L\le N\).
\end{proof}

\begin{lemma}[Spanning complementary identities give componentwise constraints]
    \label{lem:block_diagonal_boundary_component_complementary_word_identities}
    \lean{MPSTensor.blockDiagonal_boundary_component_chainGroundSpace_of_complementary_word_identities}
    \leanok
    \uses{thm:word_span_common_boundary_matrix,
        lem:block_diagonal_boundary_component_right_boundary_matrices}
    Assume \(0<N\), \(L\le N\), and that for every block \(j\), the
    length-\(N-L\) word products of \(A_j\) span the full matrix algebra. Suppose
    that, for every block \(j\), every boundary-crossing cyclic interval
    beginning at \(i\), and every complementary word
    \(\rho\) of length \(N-L\), there is a matrix \(E_{j,i,\rho}\) such
    that for every word \(\beta\) of length \(i+L-N\),
    \[
        \mu_j^N X_j A^j_\beta A^j_\rho
        =
        A^j_\beta E_{j,i,\rho}.
    \]
    Then, for every block \(j\),
    \[
        \Gamma_N^{A_j}(\mu_j^NX_j)\in\mathcal G_{N,L}(A_j).
    \]
\end{lemma}

\begin{proof}
    \leanok
    \uses{thm:word_span_common_boundary_matrix,
        lem:block_diagonal_boundary_component_right_boundary_matrices}
    For fixed \(j,i\), apply
    Theorem~\ref{thm:word_span_common_boundary_matrix} to the family
    \[
        Z_\beta=\mu_j^NX_jA^j_\beta,\qquad F_\beta=A^j_\beta.
    \]
    Since the length-\(N-L\) complementary-word products span the full matrix
    algebra, there is a single matrix \(Y_{j,i}\) such that, for every \(\beta\),
    \[
        \mu_j^NX_jA^j_\beta=A^j_\beta Y_{j,i}
    \]
    Lemma~\ref{lem:block_diagonal_boundary_component_right_boundary_matrices}
    then gives the componentwise periodic-chain constraints.
\end{proof}

\begin{lemma}[Complementary identities under block injectivity give componentwise constraints]
    \label{lem:block_diagonal_boundary_component_complementary_word_identities_injective}
    \lean{MPSTensor.blockDiagonal_boundary_component_chainGroundSpace_of_complementary_word_identities_of_injective}
    \leanok
    \uses{lem:block_diagonal_boundary_component_complementary_word_identities,
        thm:wordspan_propagation_unital}
    Assume \(0<N\), \(L\le N\), and \(L+L_0\le N\). Suppose each block \(A_j\)
    is \(L_0\)-block-injective and satisfies the normalization equation
    \[
        \sum_a A^j_aA^{j\dagger}_a=I .
    \]
    Suppose also that the boundary-crossing matrix identities of
    Lemma~\ref{lem:block_diagonal_boundary_component_complementary_word_identities}
    hold for every complementary word \(\rho\) of length \(N-L\). Then, for
    every block \(j\),
    \[
        \Gamma_N^{A_j}(\mu_j^NX_j)\in\mathcal G_{N,L}(A_j).
    \]
\end{lemma}

\begin{proof}
    \leanok
    \uses{lem:block_diagonal_boundary_component_complementary_word_identities,
        thm:wordspan_propagation_unital}
    Since \(L+L_0\le N\), the complementary length satisfies \(L_0\le N-L\).
    Theorem~\ref{thm:wordspan_propagation_unital} therefore gives, for every
    block \(j\),
    \[
        S_{N-L}(A_j)=\MN{D_j}
    \]
    Apply
    Lemma~\ref{lem:block_diagonal_boundary_component_complementary_word_identities}.
\end{proof}

\begin{lemma}[Normalization for words]
    \label{lem:sum_evalWord_mul_conjTranspose_evalWord}
    \lean{MPSTensor.sum_evalWord_mul_conjTranspose_evalWord}
    \leanok
    \uses{}
    If
    \[
        \sum_a A_aA_a^\dagger=I ,
    \]
    then, for every length \(M\),
    \[
        \sum_\rho A_\rho A_\rho^\dagger=I ,
    \]
    where \(\rho\) ranges over all words of length \(M\).
\end{lemma}

\begin{proof}
    \leanok
    \uses{}
    The case \(M=0\) is the identity matrix.  For the induction step, write a
    word of length \(M+1\) as \(a\rho\).  Since \(A_{a\rho}=A_aA_\rho\),
    \[
        \sum_{a,\rho} A_aA_\rho A_\rho^\dagger A_a^\dagger
        =
        \sum_a A_a
        \left(\sum_\rho A_\rho A_\rho^\dagger\right)
        A_a^\dagger
        =
        \sum_a A_aA_a^\dagger
        =
        I .
    \]
\end{proof}

\begin{lemma}[Compatibility hypotheses under block injectivity give componentwise constraints]
    \label{lem:block_diagonal_boundary_component_pgvwc_comparison_injective}
    \lean{MPSTensor.blockDiagonal_boundary_component_chainGroundSpace_of_pgvwc_comparison_of_injective}
    \leanok
    \uses{lem:sum_evalWord_mul_conjTranspose_evalWord,
        lem:complementary_word_boundary_identities_from_pgvwc,
        lem:block_diagonal_boundary_component_complementary_word_identities_injective}
    Assume \(0<N\), \(L\le N\), and \(L+L_0\le N\). Suppose each block \(A_j\)
    is \(L_0\)-block-injective and satisfies
    \[
        \sum_a A^j_aA^{j\dagger}_a=I .
    \]
    Fix block-diagonal boundary conditions \(X_j\). For every boundary-crossing
    interval beginning at \(i\), suppose there are matrices \(C^j_{i,\rho}\),
    indexed by complementary words \(\rho\) of length \(N-L\), such that for
    every wrapped word \(\beta\) of length \(i+L-N\),
    \[
        A^j_\beta C^j_{i,\rho}
        =
        \bigl((\mu_j^NX_j)A^j_\beta\bigr)A^j_\rho .
    \]
    Then, for every block \(j\),
    \[
        \Gamma_N^{A_j}(\mu_j^NX_j)\in\mathcal G_{N,L}(A_j).
    \]
\end{lemma}

\begin{proof}
    \leanok
    \uses{lem:sum_evalWord_mul_conjTranspose_evalWord,
        lem:complementary_word_boundary_identities_from_pgvwc,
        lem:block_diagonal_boundary_component_complementary_word_identities_injective}
    The normalization equation gives, by induction on word length,
    \[
        \sum_\rho A^j_\rho A^{j\dagger}_\rho=I
    \]
    for words \(\rho\) of length \(N-L\). Applying
    Lemma~\ref{lem:complementary_word_boundary_identities_from_pgvwc} with
    \(X=\mu_j^NX_j\) gives, for every complementary word \(\rho\), a matrix
    \(E_{j,i,\rho}\) such that for every wrapped word \(\beta\),
    \[
        \bigl((\mu_j^NX_j)A^j_\beta\bigr)A^j_\rho
        =
        A^j_\beta E_{j,i,\rho}.
    \]
    Lemma~\ref{lem:block_diagonal_boundary_component_complementary_word_identities_injective}
    then gives the componentwise periodic-chain constraint.
\end{proof}

\begin{lemma}[Trace decompositions under block injectivity give componentwise constraints]
    \label{lem:block_diagonal_boundary_component_trace_decomposition_injective}
    \lean{MPSTensor.blockDiagonal_boundary_component_chainGroundSpace_of_trace_decomposition_of_injective}
    \leanok
    \uses{lem:complementary_word_boundary_identities_from_trace_decompositions,
        lem:block_diagonal_boundary_component_complementary_word_identities_injective}
    Assume \(0<N\), \(L\le N\), and \(L+L_0\le N\). Suppose each block \(A_j\)
    is \(L_0\)-block-injective and satisfies
    \[
        \sum_a A^j_aA^{j\dagger}_a=I .
    \]
    Assume also that the simultaneous tuples \(w\mapsto(A^1_w,\ldots,A^g_w)\)
    at the middle-word length \(m\) span the full product algebra
    \(\prod_jM_{D_j}(\mathbb C)\).
    Fix block-diagonal boundary conditions \(X_j\). For every boundary-crossing
    interval beginning at \(i\), suppose there are matrices \(C^j_{i,\rho}\),
    indexed by complementary words \(\rho\) of length \(N-L\), such that for
    every wrapped word \(\beta\), complementary word \(\rho\), and middle word
    \(w\),
    \[
        \sum_j\operatorname{tr}\!\left(A^j_\beta C^j_{i,\rho}A^j_w\right)
        =
        \sum_j
        \operatorname{tr}\!\left(((\mu_j^NX_j)A^j_\beta)A^j_\rho A^j_w\right).
    \]
    Then, for every block \(j\),
    \[
        \Gamma_N^{A_j}(\mu_j^NX_j)\in\mathcal G_{N,L}(A_j).
    \]
\end{lemma}

\begin{proof}
    \leanok
    \uses{lem:complementary_word_boundary_identities_from_trace_decompositions,
        lem:block_diagonal_boundary_component_complementary_word_identities_injective}
    Applying
    Lemma~\ref{lem:complementary_word_boundary_identities_from_trace_decompositions}
    with \(X_j\) replaced by \(\mu_j^NX_j\), for every block \(j\), every
    boundary-crossing interval \(i\), and every complementary word \(\rho\),
    this gives a matrix \(E_{j,i,\rho}\) such that, for every wrapped word
    \(\beta\),
    \[
        ((\mu_j^NX_j)A^j_\beta)A^j_\rho=A^j_\beta E_{j,i,\rho}
    \].
    These are precisely the complementary-word
    identities of
    Lemma~\ref{lem:block_diagonal_boundary_component_complementary_word_identities_injective},
    which gives the periodic-chain constraint for each block.
\end{proof}

\begin{lemma}[Complementary identities give periodic block components]
    \label{lem:block_diagonal_boundary_periodic_components_from_complementary_identities}
    \lean{MPSTensor.exists_blockDiagonal_boundary_chainGroundSpace_of_complementary_identities_bnt_c1}
    \leanok
    \uses{lem:bnt_block_diagonal_open_boundary_matrices,
        lem:block_diagonal_boundary_component_complementary_word_identities_injective}
    With the hypotheses and notation of
    Lemma~\ref{lem:bnt_block_diagonal_open_boundary_matrices}, assume in
    addition that \(L+L_0\le N\).  Let
    \[
        \psi\in
        \mathcal G_{N,L}\!\left(\bigoplus_j\mu_jA_j\right).
    \]
    Suppose that whenever
    \[
        \psi=\Gamma_N^{\oplus_j\mu_jA_j}\!\left(\bigoplus_jX_j\right),
    \]
    the corresponding boundary conditions satisfy the complementary-word
    identities: for every boundary-crossing interval \(i\), every wrapped
    word \(\beta\), and every complementary word \(\rho\),
    \[
        \mu_j^N X_jA^j_\beta A^j_\rho=A^j_\beta E_{j,i,\rho}
    \]
    Then there are block-diagonal boundary conditions \(X_j\) with
    \[
        \psi=\Gamma_N^{\oplus_j\mu_jA_j}\!\left(\bigoplus_jX_j\right)
    \]
    and, for every \(j\),
    \[
        \Gamma_N^{A_j}(\mu_j^NX_j)\in\mathcal G_{N,L}(A_j)
    \]
\end{lemma}

\begin{proof}
    \leanok
    \uses{lem:bnt_block_diagonal_open_boundary_matrices,
        lem:block_diagonal_boundary_component_complementary_word_identities_injective}
    By Lemma~\ref{lem:bnt_block_diagonal_open_boundary_matrices} there are
    matrices \(X_j\) with
    \[
        \psi=\Gamma_N^{\oplus_j\mu_jA_j}\!\left(\bigoplus_jX_j\right).
    \]
    These \(X_j\) satisfy the assumed complementary-word identities, so
    Lemma~\ref{lem:block_diagonal_boundary_component_complementary_word_identities_injective}
    gives, for every \(j\),
    \[
        \Gamma_N^{A_j}(\mu_j^NX_j)\in\mathcal G_{N,L}(A_j)
    \]
\end{proof}
